\documentclass[11pt]{article}
\usepackage[margin=1in]{geometry}
\usepackage{hyperref}
\usepackage{listings}
\usepackage{xcolor}
\usepackage{longtable}
\usepackage{booktabs}
\usepackage{enumitem}
\usepackage{amsmath}

\lstset{
    basicstyle=\ttfamily\small,
    backgroundcolor=\color{gray!10},
    frame=single,
    framerule=0.5pt,
    breaklines=true,
    showstringspaces=false,
    columns=flexible
}

\title{FEMM Magnetics File Formats\\[0.5em]
\large Input (\texttt{.fem}) and Output (\texttt{.ans}) Specifications}
\author{D.\ Meeker}
\date{\today}

\begin{document}
\maketitle
\tableofcontents
\newpage

% ======================================================================
\section{Overview}

The FEMM magnetics solver (\texttt{fkn}) reads a problem definition from
a \texttt{.fem} file and writes the solution to an \texttt{.ans} file.
The \texttt{.ans} file begins with a verbatim copy of the input
\texttt{.fem} file, followed by a \texttt{[Solution]} section containing
the finite element mesh and computed field values.

All floating-point values are written with 17 significant digits
(\verb|%.17g| format) to preserve full double precision across
read/write cycles.  Fields within a line are separated by tabs or
spaces.  Section headers use square brackets (\verb|[Tag]|); property
block delimiters use angle brackets (\verb|<BeginBlock>| /
\verb|<EndBlock>|).  String comparisons during parsing are
case-insensitive.

% ======================================================================
\section{The \texttt{.fem} Input File}

% ----------------------------------------------------------------------
\subsection{Global Problem Parameters}

Each parameter occupies one line in the format
\verb|[Tag] = value|.  All are optional; defaults are applied if
missing.

\begin{longtable}{lllp{6cm}}
\toprule
Tag & Type & Default & Description \\
\midrule
\endhead
\verb|[Format]|      & double  & 4.0        & File format version \\
\verb|[Frequency]|   & double  & 0          & AC frequency (Hz); 0 = DC \\
\verb|[Precision]|   & double  & 1e-8       & Solver convergence tolerance \\
\verb|[MinAngle]|    & double  & 30         & Minimum mesh element angle (degrees) \\
\verb|[DoSmartMesh]| & int     & 1          & Adaptive meshing (0=off, 1=on) \\
\verb|[Depth]|       & double  & 1.0        & Extrusion depth (planar problems) \\
\verb|[LengthUnits]| & string  & millimeters & inches, millimeters, centimeters, meters, mils, microns \\
\verb|[ProblemType]| & string  & planar     & planar or axisymmetric \\
\verb|[Coordinates]| & string  & cartesian  & cartesian or polar \\
\verb|[ACSolver]|    & int     & 0          & 0=successive approximation, 1=Newton \\
\verb|[PrevType]|    & int     & 0          & 0=none, 1=incremental, 2=frozen, 3=restart \\
\verb|[PrevSoln]|    & string  & (empty)    & Path to previous solution (quoted) \\
\verb|[Comment]|     & string  & (empty)    & Problem description (quoted; \verb|\n| = newline) \\
\bottomrule
\end{longtable}

\subsubsection{Axisymmetric External Region}

Written only when \verb|ProblemType| is axisymmetric and both
\verb|extRo| and \verb|extRi| are nonzero:

\begin{longtable}{llp{8cm}}
\toprule
Tag & Type & Description \\
\midrule
\verb|[extZo]| & double & Outer Z coordinate \\
\verb|[extRo]| & double & Outer radius \\
\verb|[extRi]| & double & Inner radius \\
\bottomrule
\end{longtable}

% ----------------------------------------------------------------------
\subsection{Point (Nodal) Properties}

\begin{lstlisting}
[PointProps] = <count>
<BeginPoint>
  <PointName> = "name"
  <I_re> = 0        real part of point current (A)
  <I_im> = 0        imaginary part (AC only)
  <A_re> = 0        real part of prescribed A
  <A_im> = 0        imaginary part (AC only)
<EndPoint>
... repeated count times
\end{lstlisting}

% ----------------------------------------------------------------------
\subsection{Boundary Properties}

\begin{lstlisting}
[BdryProps] = <count>
<BeginBdry>
  <BdryName> = "name"
  <BdryType> = 0
  <A_0> = 0          <A_1> = 0          <A_2> = 0
  <Phi> = 0
  <c0> = 0           <c0i> = 0
  <c1> = 0           <c1i> = 0
  <Mu_ssd> = 0       <Sigma_ssd> = 0
  <innerangle> = 0   <outerangle> = 0
<EndBdry>
\end{lstlisting}

Boundary types:
\begin{itemize}[nosep]
  \item 0 = prescribed A
  \item 1 = small skin depth
  \item 2 = mixed
  \item 3 = strategic dual image (SDI)
  \item 4 = periodic
  \item 5 = antiperiodic
  \item 6 = periodic air gap element (AGE)
  \item 7 = antiperiodic air gap element (AGE)
\end{itemize}

% ----------------------------------------------------------------------
\subsection{Material (Block) Properties}

\begin{lstlisting}
[BlockProps] = <count>
<BeginBlock>
  <BlockName> = "name"
  <Mu_x> = 1         relative permeability, x-direction
  <Mu_y> = 1         relative permeability, y-direction
  <H_c> = 0          coercivity (A/m)
  <H_cAngle> = 0     magnetization direction (degrees)
  <J_re> = 0         source current density, real (MA/m^2)
  <J_im> = 0         source current density, imag (MA/m^2)
  <Sigma> = 0        conductivity (S/m)
  <d_lam> = 0        lamination thickness (mm)
  <Phi_h> = 0        hysteresis lag angle (degrees)
  <Phi_hx> = 0       hysteresis lag angle, x (degrees)
  <Phi_hy> = 0       hysteresis lag angle, y (degrees)
  <LamType> = 0      lamination type
  <LamFill> = 1      lamination fill factor (0--1)
  <NStrands> = 0     number of wire strands
  <WireD> = 0        wire diameter (mm)
  <BHPoints> = N
    B_0  H_0         B-H curve data (Tesla, A/m)
    B_1  H_1         ... repeated N times
    ...
<EndBlock>
\end{lstlisting}

Lamination types:
\begin{itemize}[nosep]
  \item 0 = not laminated
  \item 1 = laminated in x (parallel to x)
  \item 2 = laminated in y (parallel to y)
  \item 3 = magnet wire
  \item 4 = plain stranded
  \item 5 = litz wire
  \item 6 = rectangular wire
  \item 7 = 10\% CCA
  \item 8 = 15\% CCA
\end{itemize}

% ----------------------------------------------------------------------
\subsection{Circuit Properties}

\begin{lstlisting}
[CircuitProps] = <count>
<BeginCircuit>
  <CircuitName> = "name"
  <TotalAmps_re> = 0     total current, real (A)
  <TotalAmps_im> = 0     total current, imag (A)
  <CircuitType> = 0      0=series, 1=parallel
<EndCircuit>
\end{lstlisting}

% ----------------------------------------------------------------------
\subsection{Geometry}

\subsubsection{Nodes}

\begin{lstlisting}
[NumPoints] = <count>
x    y    propIdx    inGroup    [lfs]
\end{lstlisting}

\begin{itemize}[nosep]
  \item \texttt{x}, \texttt{y}: coordinates (double)
  \item \texttt{propIdx}: index into point property list (1-based; 0 = none)
  \item \texttt{inGroup}: group number (int)
  \item \texttt{lfs}: local feature size constraint (double, optional; $-1$ = none)
\end{itemize}

\subsubsection{Line Segments}

\begin{lstlisting}
[NumSegments] = <count>
n0    n1    maxSideLength    propIdx    hidden    inGroup
\end{lstlisting}

\begin{itemize}[nosep]
  \item \texttt{n0}, \texttt{n1}: endpoint node indices (0-based)
  \item \texttt{maxSideLength}: mesh size constraint (double; $-1$ = auto)
  \item \texttt{propIdx}: boundary property index (1-based; 0 = none)
  \item \texttt{hidden}: 0 = visible, 1 = suppressed from mesh
  \item \texttt{inGroup}: group number
\end{itemize}

\subsubsection{Arc Segments}

\begin{lstlisting}
[NumArcSegments] = <count>
n0  n1  arcAngle  maxSideLength  propIdx  hidden  inGroup  mySideLength
\end{lstlisting}

\begin{itemize}[nosep]
  \item \texttt{n0}, \texttt{n1}: endpoint node indices (0-based)
  \item \texttt{arcAngle}: arc span (degrees)
  \item \texttt{maxSideLength}: mesh constraint (double)
  \item \texttt{propIdx}: boundary property index (1-based; 0 = none)
  \item \texttt{hidden}: 0 = visible, 1 = suppressed
  \item \texttt{inGroup}: group number
  \item \texttt{mySideLength}: computed side length (0 = auto)
\end{itemize}

\subsubsection{Holes}

\begin{lstlisting}
[NumHoles] = <count>
x    y    inGroup
\end{lstlisting}

Holes are seed points for regions excluded from the mesh.  They
correspond to block labels with material \texttt{"<No Mesh>"}.

\subsubsection{Block Labels}

\begin{lstlisting}
[NumBlockLabels] = <count>
x  y  blockIdx  meshSize  circIdx  magDir  inGroup  turns  extFlags  ["magDirFctn"]
\end{lstlisting}

\begin{itemize}[nosep]
  \item \texttt{x}, \texttt{y}: label position (double)
  \item \texttt{blockIdx}: material index (1-based; 0 = unassigned)
  \item \texttt{meshSize}: element size (double; $\leq 0$ = auto).
        Converted to area: $A = \pi r^2/4$
  \item \texttt{circIdx}: circuit index (1-based; 0 = none)
  \item \texttt{magDir}: magnetization direction (degrees)
  \item \texttt{inGroup}: group number
  \item \texttt{turns}: number of turns
  \item \texttt{extFlags}: bit 0 = external region, bit 1 = default region
  \item \texttt{magDirFctn}: magnetization direction function (optional, quoted string)
\end{itemize}


% ======================================================================
\section{The \texttt{.ans} Output File}

The \texttt{.ans} file begins with a verbatim copy of the input
\texttt{.fem} file (all sections described in Section~2), followed by
the solver's solution data.

% ----------------------------------------------------------------------
\subsection{Solution Header}

Immediately after the echoed \texttt{.fem} content:

\begin{lstlisting}
[Solution]
\end{lstlisting}

% ----------------------------------------------------------------------
\subsection{Mesh Nodes}

\begin{lstlisting}
<NumNodes>
\end{lstlisting}

followed by \texttt{NumNodes} lines.  The format depends on the problem
type:

\paragraph{DC (Frequency = 0), no previous solution:}
\begin{lstlisting}
x    y    A_re
\end{lstlisting}

\paragraph{DC with incremental/frozen permeability:}
\begin{lstlisting}
x    y    A_re    bc    A_prev
\end{lstlisting}

\paragraph{AC (Frequency $> 0$), no previous solution:}
\begin{lstlisting}
x    y    A_re    A_im
\end{lstlisting}

\paragraph{AC with incremental/frozen permeability:}
\begin{lstlisting}
x    y    A_re    A_im    bc    A_prev
\end{lstlisting}

\begin{itemize}[nosep]
  \item \texttt{x}, \texttt{y}: node coordinates in model units
  \item \texttt{A\_re}, \texttt{A\_im}: magnetic vector potential (Wb/m)
  \item \texttt{bc}: boundary condition index ($-1$ = none)
  \item \texttt{A\_prev}: vector potential from previous solution
\end{itemize}

Node indices are 0-based throughout the solution section.

% ----------------------------------------------------------------------
\subsection{Mesh Elements}

\begin{lstlisting}
<NumElements>
\end{lstlisting}

followed by \texttt{NumElements} lines:

\begin{lstlisting}
p0    p1    p2    labelIdx    e0    e1    e2    J_prev
\end{lstlisting}

\begin{itemize}[nosep]
  \item \texttt{p0}, \texttt{p1}, \texttt{p2}: vertex node indices (0-based)
  \item \texttt{labelIdx}: block label index (0-based, into the block label
        list from the \texttt{.fem} header)
  \item \texttt{e0}, \texttt{e1}, \texttt{e2}: boundary property index applied
        to each of the element's three edges, or $-1$ where the edge carries
        no boundary condition
  \item \texttt{J\_prev}: current density carried over from the previous
        solution, in \textbf{MA/m$^2$} --- the solver's working unit, not the
        A/m$^2$ that the postprocessor reports.  Consumers must scale by
        $10^6$ to combine it with current densities obtained from the
        postprocessor.
\end{itemize}

The static solver always writes all eight fields.  The harmonic solver
writes the first seven and appends \texttt{J\_prev} only for an incremental
(frozen-permeability) run.

\medskip\noindent
\textbf{Note on \texttt{J\_prev}'s position.}  It is the \emph{eighth} field,
after the three edge indices.  Original FEMM's postprocessor reads it as the
fifth --- picking up \texttt{e0}, a boundary index that is $-1$ on every
interior element --- which corrupts its 1$\times$ incremental Lorentz force
and torque integrals (block integral types 28--30).  femm-qt reads the field
the solver actually writes, so those three integrals differ from original
FEMM by design; every other quantity is unaffected.

Elements are written in Cuthill-McKee order (sorted by sum of node
indices), which clusters spatially nearby elements at nearby indices for
efficient postprocessor lookup.

% ----------------------------------------------------------------------
\subsection{Circuit Solution Data}

\begin{lstlisting}
<NumCircuits>
\end{lstlisting}

followed by \texttt{NumCircuits} lines (one per block label in the
problem):

\paragraph{DC:}
\begin{lstlisting}
caseType    value_re
\end{lstlisting}

\paragraph{AC:}
\begin{lstlisting}
caseType    value_re    value_im
\end{lstlisting}

\begin{itemize}[nosep]
  \item \texttt{caseType}:
    \begin{itemize}[nosep]
      \item 0 = voltage gradient (value = $\Delta V$, voltage drop)
      \item 1 = current density (value = $J$)
      \item 2 = flux constraint (from solver's flux equation)
    \end{itemize}
  \item \texttt{value\_re}, \texttt{value\_im}: solution values for the
        circuit associated with this block label
\end{itemize}

% ----------------------------------------------------------------------
\subsection{Periodic Boundary Conditions}

\begin{lstlisting}
<NumPBCs>
\end{lstlisting}

followed by \texttt{NumPBCs} lines:

\begin{lstlisting}
nodeIdx0    nodeIdx1    type
\end{lstlisting}

\begin{itemize}[nosep]
  \item \texttt{nodeIdx0}, \texttt{nodeIdx1}: paired node indices (0-based)
  \item \texttt{type}: 0 = periodic, 1 = antiperiodic
\end{itemize}

% ----------------------------------------------------------------------
\subsection{Air Gap Elements}

\begin{lstlisting}
<NumAGEs>
\end{lstlisting}

For each air gap element:

\paragraph{Line 1:} Boundary name (quoted string)
\begin{lstlisting}
"BoundaryName"
\end{lstlisting}

\paragraph{Line 2:} Parameters
\begin{lstlisting}
bdryFmt  innerAngle  outerAngle  ri  ro  totalArc  agc_re  agc_im  numArcElems  innerShift  outerShift
\end{lstlisting}

\begin{itemize}[nosep]
  \item \texttt{bdryFmt}: 0 = periodic, 1 = antiperiodic
  \item \texttt{innerAngle}, \texttt{outerAngle}: boundary angles (degrees)
  \item \texttt{ri}, \texttt{ro}: inner and outer radii (model units)
  \item \texttt{totalArc}: total arc length (degrees)
  \item \texttt{agc\_re}, \texttt{agc\_im}: air gap coefficient (complex)
  \item \texttt{numArcElems}: number of arc elements
  \item \texttt{innerShift}, \texttt{outerShift}: boundary shifts
\end{itemize}

\paragraph{Lines 3 to $2 + \text{numArcElems} + 1$:} Quadrature points
\begin{lstlisting}
n0  w0  n1  w1  n2  w2  n3  w3
\end{lstlisting}

\begin{itemize}[nosep]
  \item \texttt{n0}--\texttt{n3}: node indices for interpolation
  \item \texttt{w0}--\texttt{w3}: interpolation weights
\end{itemize}

% ----------------------------------------------------------------------
\subsection{Air Gap Ring Node Lists (optional)}

The quadrature-point mapping above is sorted at the rotor angle of the
solution that wrote it.  When air gap elements are present, the solver
appends a trailing section carrying each element's raw ring node lists,
which are independent of the rotor angle:

\begin{lstlisting}
[AgeRings]
\end{lstlisting}

For each air gap element, a count line equal to \texttt{numArcElems},
followed by that many lines:

\begin{lstlisting}
innerNodeIdx  outerNodeIdx
\end{lstlisting}

listing the mesh nodes (0-based) on the inner and outer gap rings of the
meshed sector.  A transient analysis uses these to reuse the previous
solution's mesh while rebuilding the quadrature-point pairing at the
current step's rotor angle.  The section is strictly trailing so that
readers which parse the air gap elements positionally and stop are
unaffected; when it is absent, the solver falls back to remeshing.

% ======================================================================
\section{Index Conventions}

\begin{itemize}
  \item \textbf{Property indices} (materials, boundaries, circuits, point
        properties) are \textbf{1-based} in the file.  An index of 0
        means ``no property assigned.''  The parser converts to 0-based
        internally.

  \item \textbf{Node indices} (in segments, arcs, and the solution
        section) are \textbf{0-based}.

  \item \textbf{Block label indices} in the solution element list are
        \textbf{0-based}, referencing the block label array from the
        \texttt{.fem} header.
\end{itemize}

% ======================================================================
\section{Unit Conversion}

The \verb|[LengthUnits]| tag determines the coordinate system.  All
coordinates in the file are in the specified units.  Conversion to
meters:

\begin{center}
\begin{tabular}{ll}
\toprule
Unit & Meters per unit \\
\midrule
inches      & 0.0254 \\
millimeters & 0.001 \\
centimeters & 0.01 \\
meters      & 1.0 \\
mils        & 2.54$\times 10^{-5}$ \\
microns     & $10^{-6}$ \\
\bottomrule
\end{tabular}
\end{center}

\end{document}
